% !TEX program = xelatex
%% =====================================================================
%%  MATURITNÍ PRÁCE – Gymnázium Jírovcova 8, České Budějovice
%%  Autoři šablony: Jonáš Havelka, Michal Kočer, Daniel Sýkora (MIT)
%% =====================================================================
%%  PŘEKLÁDEJTE XeLaTeXem!
%%    Overleaf: Menu → Compiler: XeLaTeX, Main document: mprace.tex
%%    Při pdfLaTeXu skončí překlad chybou (fontspec / polyglossia).
%%
%%  Co je kde:
%%    mprace.tex        tento soubor – údaje o práci, úvodní strany,
%%                      pořadí kapitol a seznamy
%%    kapitoly/*.tex    text práce, každá část ve vlastním souboru
%%    Bibliografie.bib  seznam literatury
%%    images/           obrázky
%%    vzor.pdf          vzorová práce s ukázkami (citace, rovnice, obrázky,
%%                      tabulky, diagramy, výpisy kódu…); zdrojový kód
%%                      ukázek je ve vzor/ukazky/
%%    MP.sty            styl práce – neupravujte
%% =====================================================================
\documentclass[12pt]{report}
\usepackage{MP}								% Styl maturitní práce
\addbibresource{Bibliografie.bib}			% Soubor se seznamem literatury

%% Volitelné balíčky – odkomentujte, jen pokud je potřebujete (prodlužují překlad):
% \usepackage{tikz}\usetikzlibrary{arrows.meta,shapes.geometric,positioning} % diagramy, schémata
% \usepackage[version=4]{mhchem}		% chemické vzorce \ce{...}

%% ---------------------------------------------------------------------
%%  ÚDAJE O PRÁCI – vyplňte
%% ---------------------------------------------------------------------
\author{Jméno Příjmení}						% Vaše jméno a příjmení
\title{Název maturitní práce}				% Název práce
\date{datum odevzdání}						% Např. 14. února 2027
\vedouci{Mgr. Jméno Příjmení}				% Vedoucí práce i s tituly
\place{V Českých Budějovicích}
\skolnirok{2026/2027}
\logo{\includegraphics[scale=1.25]{GJ8_logotyp}}

%% Rychlejší překlad delší práce: odkomentujte a uveďte jen kapitoly,
%% na kterých právě pracujete. Ostatní kapitoly v PDF dočasně chybí, ale
%% čísla stran a odkazy na ně zůstanou z posledního úplného překladu.
%% Před odevzdáním řádek znovu zakomentujte!
% \includeonly{kapitoly/prakticka}

\begin{document}
%% ---- ÚVODNÍ STRANY (číslováno římsky, čísla stran se nezobrazují) ----
\pagenumbering{roman}
\mytitlepage

\prohlaseni{
	Prohlašuji, že jsem tuto práci vypracoval samostatně s vyznačením všech použitých pramenů.

	% Pokud jste při psaní práce využili nástroje umělé inteligence, uveďte to zde
	% (který nástroj a k čemu). Upravte nebo smažte následující větu podle skutečnosti:
	Pro kontrolu pravopisu a úpravu formulací jsem využil nástroj umělé inteligence (název nástroje).
}

\abstrakt{
	% Shrnutí práce (10–15 vět): co jste řešili, jakou metodou a co jste zjistili.
	Text abstraktu.
}{
	% Klíčová slova oddělená čárkami
	klíčové slovo, klíčové slovo, klíčové slovo
}

\podekovani{
	% Poděkování vedoucímu práce, rodině, institucím apod.
	Děkuji vedoucímu práce \vedouci{} za odborné vedení a cenné rady.
}

{\tableofcontents\newpage}

%% ---- VLASTNÍ PRÁCE (číslováno arabsky) ----
\pagenumbering{arabic}

% !TEX root = ../mprace.tex
%% ÚVOD (zhruba 1–2 strany)
%%  - téma práce a proč jste si ho vybrali,
%%  - cíle práce (co chcete zjistit, vytvořit nebo ověřit),
%%  - stručně, co čtenář najde v jednotlivých kapitolách.
\chapter*{Úvod}
\addcontentsline{toc}{chapter}{Úvod}

Text úvodu.


\part{Teoretická část}						% Přejmenujte podle obsahu práce
% !TEX root = ../mprace.tex
%% TEORETICKÁ ČÁST – z čeho vycházíte: pojmy, principy, přehled literatury.
%% Kapitol může být víc; každá začíná příkazem \chapter, podkapitola \section.
%%
%% Ukázky najdete ve vzor.pdf, zdrojový kód ve složce vzor/ukazky/:
%%   citace.tex      citace, české uvozovky, tabulka
%%   zkratky.tex     seznam zkratek (nepovinný)
%%   matematika.tex  rovnice, odkazy na rovnice
%%   obrazky.tex     obrázky, obrázky vedle sebe, obtékání textem
\chapter{Název první kapitoly}

\section{Název podkapitoly}

Text podkapitoly.

\chapter{Název druhé kapitoly}

Text kapitoly.


\part{Praktická část}						% Přejmenujte podle obsahu práce
% !TEX root = ../mprace.tex
%% PRAKTICKÁ ČÁST – co jste udělali: postup, měření, program, výsledky.
%% Kapitol může být víc; každá začíná příkazem \chapter, podkapitola \section.
%%
%% Ukázky najdete ve vzor.pdf, zdrojový kód ve složce vzor/ukazky/:
%%   jednotky.tex    jednotky SI (siunitx)
%%   chemie.tex      chemické vzorce a rovnice (odkomentujte mhchem v mprace.tex)
%%   tikz.tex        nákresy a vývojové diagramy (odkomentujte tikz v mprace.tex)
%%   schemata.tex    bloková schémata, UML, schéma databáze (také tikz)
%%   vypisy.tex      výpisy programů
\chapter{Název kapitoly praktické části}

\section{Název podkapitoly}

Text podkapitoly.


% !TEX root = ../mprace.tex
%% ZÁVĚR – shrnutí výsledků, zhodnocení splnění cílů z úvodu,
%% případně návrhy, jak v práci pokračovat.
\chapter*{Závěr}
\addcontentsline{toc}{chapter}{Závěr}

Text závěru.


%% ---- SEZNAMY (vytvářejí se automaticky; nepotřebné smažte) ----
\printbibliography[heading=bibintoc]		% Seznam literatury (jen citované zdroje)
\listoffigures								% Seznam obrázků
\listoftables								% Seznam tabulek
% \lstlistoflistings						% Seznam výpisů programů

%% ---- PŘÍLOHY (nemáte-li přílohy, smažte celé prostředí appendices) ----
\begin{appendices}
% !TEX root = ../mprace.tex
%% PŘÍLOHY – fotografie, delší tabulky a výpisy, seznam souborů v archivu…
%% Každá příloha začíná příkazem \chapter (číslují se A, B, C…).
%% Ukázka: vzor/ukazky/prilohy.tex
%% Nemáte-li přílohy, smažte v mprace.tex celé prostředí appendices.
\chapter{Název přílohy}

Obsah přílohy.

\end{appendices}

\end{document}
